\documentclass[12pt,a4paper]{article}
\usepackage[utf8]{inputenc}
\usepackage[portuguese]{babel}
\usepackage{amsmath}
\usepackage{graphicx}
\usepackage{hyperref}
\usepackage{geometry}

\geometry{a4paper, margin=2cm}

\title{Relatorio de Desenvolvimento: Algoritmo de Deteccao e Prevencao de Quedas para Iniciacao Cientifica}
\author{Inovacao UFLA}
\date{\today}

\begin{document}

\maketitle

\section{Introducao}
Este documento resume o desenvolvimento do algoritmo de visao computacional voltado para a deteccao e analise de risco de quedas em pacientes. O projeto visa oferecer uma solucao em tempo real, utilizando uma webcam para monitorar a postura e emitir alertas automatizados caso seja identificada uma queda ou uma instabilidade severa. A ferramenta serve de apoio para registro de dados e estudos de ocorrencia na Iniciacao Cientifica.

\section{Tecnologias Utilizadas}
\begin{itemize}
    \item \textbf{Python 3:} Linguagem principal do projeto.
    \item \textbf{OpenCV (cv2):} Utilizado para captura e processamento das imagens geradas por webcams (nativas ou conectadas via USB).
    \item \textbf{MediaPipe (Google):} Ferramenta central para deteccao de \textit{landmarks} corporais (pose estimation).
    \item \textbf{NumPy, Time, Csv, Datetime:} Bibliotecas auxiliares para calculo matematico, controle de concorrencia e geracao de tabelas de registros de dados.
\end{itemize}

\section{Arquitetura e Logica de Deteccao}
\subsection{Captura e Estimativa de Pose}
Em cada quadro de video, o OpenCV coleta a imagem e a fornece a solucao \textit{MediaPipe Pose}. O algoritmo consolida dinamicamente os \textit{landmarks}, permitindo mapear nao apenas quedas ja ocorridas, mas predizer instabilidades corporais previas.

\subsection{Identificacao de Queda e Margem de Erro (Robustez)}
Calcula-se a distancia vertical instantanea ($\Delta Y$) entre o quadril e os ombros. Uma possivel queda e sinalizada quando a posicao do ombro indica inversao perante o quadril, ou quando a distancia vertical e drasticamente reduzida na base da imagem. Aplicou-se um limiar de frames (FALL\_FRAME\_THRESHOLD = 5) como Margem de Erro. Dessa forma, a metrica de queda exige sustentacao do estado da queda por 5 frames consecutivos para evitar Falsos Positivos. Foi implementado tambem um periodo de cooldown (3.0s) a fim de impedir re-checagens abusivas enquanto o individuo ja esta acidentado.

\subsection{Alerta Preditivo de Risco (Instabilidade e Oscilacao)}
Para enriquecer a abordagem profilatica exigida na literatura medica, o algoritmo foi expandido para calcular a variancia horizontal (Eixo X). O monitoramento salva a posicao dos ombros nos ultimos 15 \textit{frames}. Se a oscilacao entre os limites esquerdo e direito ultrapassar um limiar pre-estabelecido sem que o individuo se desloque verticalmente, uma tontura ou espasmo postural e interpretada, ativando antecipadamente na interface um alerta amarelo de: \textbf{"RISCO: INSTABILIDADE CORPORAL!"}.

\section{Provas de Conceito (Amostragem e Matriz de Confusao)}
A fim de submeter o algoritmo a validacoes academicas estritas de Iniciação Científica, um modulo extra secundario (\texttt{dataset\_evaluator.py}) foi acoplado ao modelo. O avaliador consome pastas de datasets (videos com quedas reais e videos de acoes de rotina - ADLs). 
Ao ser executado em lote, o roteiro processa todo o banco de imagens, discriminando estatisticamente as execucoes com uma \textbf{Matriz de Confusao}, separando: Verdadeiros Positivos, Verdadeiros Negativos, Falsos Positivos e Falsos Negativos, que por consequencia, fornecem relaboracao de Acuracia, Sensibilidade e Especificidade da rede.

\section{Registro de Dados e Tabulacao (CSV)}
Sempre que uma queda ou amostragem e confirmada em lote ou em tempo-real, os dados sao rigorosamente documentados em sistemas em formato .csv, registrando a \textit{timeline}, ID do incidente e as variaveis Y de compressao identificadas pelo \textit{Machine Learning}.

\section{Conclusao}
O algoritmo apresenta agora multiplas frentes sistemicas. O processamento unificado heuristico preve nao reativamente mas ativamente as quedas iminentes em pacientes. O sistema conta ainda com uma integracao modular de \textit{Dataset Evaluator} pronto para tabular matrizes de confusao com dados academicos extraidos de bancos de videos reais de quedas em pacientes.

\end{document}
